\section{Endpoint ledger and the annular resonance budget}
\label{sec:tpc48-endpoint}

We now insert the inherited TPC endpoint.  Write
\begin{equation}
 \mu:=\frac{267}{400},
 \qquad
 \nu:=\frac{133}{400},
 \qquad
 \kappa:=\frac1{400}.
 \label{eq:tpc48-endpoint-exponents}
\end{equation}
Thus
\begin{equation}
 M=X^{\mu+o(1)},
 \qquad
 J=X^{\nu+o(1)},
 \qquad
 \mu+\nu=1.
 \label{eq:tpc48-endpoint-scales}
\end{equation}
The actual slope support only obeys \(K\le X^{\mu+o(1)}\).  Statements
that identify \(K\) with \(M\) below are therefore model sharpness
statements for a dense slope box, not lower bounds for the literal
coefficient.

\subsection{Critical bands and exact fine tubes}

Take
\begin{equation}
 L=X^{\nu+o(1)}\asymp J.
 \label{eq:tpc48-critical-band-count}
\end{equation}
For a full scalar orbit window of length \(N\asymp J\),
\cref{thm:tpc48-band-trace} gives
\begin{equation}
 \operatorname{tr}(E_{I,v}^*\Pi_kE_{I,v})=O(1).
 \label{eq:tpc48-critical-trace}
\end{equation}
For \(m\asymp M\), \cref{cor:tpc48-tube-geometry} gives the exact
branch half-width
\begin{equation}
 \frac1{LH_0m}=X^{-1+o(1)}.
 \label{eq:tpc48-endpoint-tube-width}
\end{equation}
The fine scale is therefore recovered simultaneously across an exact
orthogonal resolution of the whole orbit-frequency circle.  This is a
geometric resolution statement; it is not a mass lower bound or a
spectral gap for the complete actual vector.

It is useful to compare this choice with the canonical band in
TPC--47.  There the band count needed to cover the quotient was
\begin{equation}
 L_{47}=X^{\nu-\kappa+o(1)}=X^{33/100+o(1)},
 \label{eq:tpc48-tpc47-cover-count}
\end{equation}
the per-band trace was \(X^{\kappa+o(1)}\), and the resonance
half-width was \(X^{-1+\kappa+o(1)}=X^{-399/400+o(1)}\).
Passing to the critical tiling adds the missing factor \(X^\kappa\)
to the band count, makes the trace bounded, and sharpens the tube
scale by the same factor.  Unlike the single projected sector in
TPC--47, the present \(L\) bands sum exactly to the identity.

\subsection{Pointwise collision ledger}

On the annular shift range below, the no-wrap hypothesis holds and
the large-sieve factor is
\begin{equation}
 D_J(t;K)=\min\left\{K,1+\frac1{Jt}\right\}.
 \label{eq:tpc48-endpoint-D}
\end{equation}
At the fine edge \(t=X^{-1+\epsilon}\), the universal upper bound
\(K\le X^{\mu+o(1)}\) gives
\begin{equation}
 D_J(t;K)\le X^{\mu-\epsilon+o(1)},
 \qquad
 \frac J{H_0}D_J(t;K)\le X^{1-\epsilon+o(1)}.
 \label{eq:tpc48-fine-edge-ledger}
\end{equation}
At the coarse edge \(t=X^{-\mu-\epsilon}\),
\begin{equation}
 D_J(t;K)\le X^{\mu-\nu+\epsilon+o(1)}
 =X^{67/200+\epsilon+o(1)},
 \label{eq:tpc48-coarse-edge-D}
\end{equation}
and hence
\begin{equation}
 \frac J{H_0}D_J(t;K)\le X^{\mu+\epsilon+o(1)}.
 \label{eq:tpc48-coarse-edge-prefactor}
\end{equation}
The corresponding amplitude prefactors are

\begin{center}
\begin{tabular}{c|c|c}
shift scale & squared prefactor & amplitude prefactor\\ \hline
\(t=X^{-1+\epsilon}\)
 & \(X^{1-\epsilon+o(1)}\)
 & \(X^{1/2-\epsilon/2+o(1)}\)\\
\(t=X^{-\mu-\epsilon}\)
 & \(X^{\mu+\epsilon+o(1)}\)
 & \(X^{\mu/2+\epsilon/2+o(1)}\)
\end{tabular}
\end{center}

For a dense slope box at the coarse scale \(t\asymp M^{-1}\), inside
the no-wrap slab \eqref{eq:tpc48-coarse-edge-slab}, the normalized
collision load is
\begin{equation}
 \frac MJ=X^{\mu-\nu+o(1)}=X^{67/200+o(1)}.
 \label{eq:tpc48-critical-collision-load}
\end{equation}
The inherited TPC allowance is the normalized determinant-packet loss
\begin{equation}
 K_B=X^{\kappa+o(1)}
 \label{eq:tpc48-inherited-allowance}
\end{equation}
from the TPC--45/47 closure ledger
\citep{WangTPC45,WangTPC47}.  The unresolved dense-model exponent gap
is therefore
\begin{equation}
 (\mu-\nu)-\kappa
 =\frac{67}{200}-\frac1{400}
 =\frac{133}{400}
 =\nu.
 \label{eq:tpc48-endpoint-allowance-gap}
\end{equation}
Thus the common-profile theorem saves the full orbit length relative
to the unrestricted \(K\)-loss, but the coarse-edge bound still does
not fit the atomic TPC budget.  Attainment of this common-profile
operator scale requires the consecutive slope cluster in
\cref{prop:tpc48-large-sieve-sharpness}; it is not asserted for the
literal coefficient.

\subsection{A logarithmic relative annulus theorem}

For
\begin{equation}
 0<\epsilon<\frac\nu2=\frac{133}{800},
 \qquad
 \cA_\epsilon
 :=\left\{\alpha\in\T:
 X^{-1+\epsilon}\le\|\alpha\|_\T
 \le X^{-\mu-\epsilon}\right\},
 \label{eq:tpc48-annulus}
\end{equation}
the annulus is nonempty for large \(X\).  Moreover,
\begin{equation}
 H_0(M_+-M_-)\|\alpha\|_\T=o(1)
 \label{eq:tpc48-annulus-no-wrap}
\end{equation}
uniformly on \(\cA_\epsilon\), after the inherited fixed dyadic
localization.  Hence \cref{thm:tpc48-actual-reduction} applies there.

Fix one global admissible signed representation \(\mathscr R\).  It
is important here not to integrate the pointwise infimum
\(\mathfrak S_{\rm pr}\), for which no measurability assertion was
made.

\begin{theorem}[Relative annular resonance budget]
\label{thm:tpc48-annular-budget}
For every fixed global admissible representation \(\mathscr R\), with
the measurable envelope specified in
\eqref{eq:tpc48-representation-envelope},
\begin{equation}
 \int_{\cA_\epsilon}
 \|\widehat{\cF}^{\rm fr,\mathrm L}(\alpha)\|^2\,d\alpha
 \le
 \frac J{H_0}
 \int_{\cA_\epsilon}
 D_J(\|\alpha\|_\T;K)
 \mathfrak S_{\mathscr R}(\alpha)^2\,d\alpha.
 \label{eq:tpc48-weighted-annular-budget}
\end{equation}
In particular, if
\begin{equation}
 \sup_{\alpha\in\cA_\epsilon}
 \mathfrak S_{\mathscr R}(\alpha)\le S_*,
 \label{eq:tpc48-annular-envelope-bound}
\end{equation}
then
\begin{equation}
 \int_{\cA_\epsilon}
 \|\widehat{\cF}^{\rm fr,\mathrm L}(\alpha)\|^2\,d\alpha
 \ll_{\epsilon,H_0} (\log X)S_*^2.
 \label{eq:tpc48-log-annular-budget}
\end{equation}
The contribution of this annulus to Fourier inversion at \(h_0\)
satisfies
\begin{equation}
 \left\|
 \int_{\cA_\epsilon}
 \widehat{\cF}^{\rm fr,\mathrm L}(\alpha)
 \e(h_0\alpha)\,d\alpha
 \right\|
 \ll
 X^{-(\mu+\epsilon)/2+o(1)}
 (\log X)^{1/2}S_*.
 \label{eq:tpc48-annular-fixed-shift-contribution}
\end{equation}
\end{theorem}

\begin{proof}
The first estimate is the squared version of
\eqref{eq:tpc48-actual-reduction}, retained for the fixed
representation \(\mathscr R\).  The elementary inequality
\begin{equation}
 D_J(t;K)\le 1+\frac1{Jt}
 \label{eq:tpc48-annulus-D-bound}
\end{equation}
holds throughout the annulus.  By symmetry around zero,
\begin{align}
 \int_{\cA_\epsilon}D_J(\|\alpha\|_\T;K)\,d\alpha
 &\ll X^{-\mu-\epsilon}
   +\frac1J\log\frac{X^{-\mu-\epsilon}}
                         {X^{-1+\epsilon}} \\
 &\ll X^{-\nu+o(1)}.
 \label{eq:tpc48-integrated-D}
\end{align}
Multiplication by \(J/H_0\) proves
\eqref{eq:tpc48-log-annular-budget}.  Finally
\(|\cA_\epsilon|\ll X^{-\mu-\epsilon}\), so Cauchy--Schwarz in
\(\alpha\) gives
\eqref{eq:tpc48-annular-fixed-shift-contribution}.
\end{proof}

The theorem is deliberately relative to \(S_*\).  Neither
\eqref{eq:tpc48-log-annular-budget} nor
\eqref{eq:tpc48-annular-fixed-shift-contribution} estimates the
projective envelope by the physical atomic diagonal.  It also leaves
the ultra-low interval, the complementary shift frequencies, and any
sparse alias reconstruction untreated.
